% GENERATED FILE. Edit canonical lessons, then run npm run generate:books.
% canonical-source-hash: fnv1a64:31b6c480ce6fab2a
% canonical-lessons: SA-C386-siddhih, SA-C386-sidhyati, SA-C386-arjayati, SA-C386-abhyudayah, SA-C386-samrddhih, SA-R386-review-380-386

\chapter{Success}
\label{ch:sa-success}
\hlchaptermodality{386}

\begin{chapteropening}
\textbf{By the end of this chapter:} \emph{I can say what will come if a plan succeeds.}

Complete the last lesson of chapter 386: I can say what will come if a plan succeeds.
\end{chapteropening}

\section[\texorpdfstring{\sk{सिद्धिः} (siddhiḥ)}{सिद्धिः (siddhiḥ)}]{\sk{सिद्धिः} (siddhiḥ) — an accomplishment}

\label{lesson:SA-C386-siddhih}



\begin{quote}
Before the new one: say the Sanskrit for rewards (with a prize), then the Sanskrit for excellence.
\end{quote}

\subsection*{You'll want to know: \sk{सिद्धिः}}
\textbf{\sk{सिद्धिः}} — \emph{siddhiḥ} — \textquotedblleft{}an accomplishment\textquotedblright{}.

\begin{quote}
\sk{यदि} \sk{प्रयत्नः} \sk{दृढः}, \sk{तर्हि} \sk{सिद्धिः} \sk{अवश्यं} \sk{भविष्यति।}
\end{quote}

\begin{quote}
\emph{yadi prayatnaḥ dṛḍhaḥ, tarhi siddhiḥ avaśyaṃ bhaviṣyati.}
\end{quote}

\begin{quote}
If the effort is firm, success will surely come.
\end{quote}

A feminine noun.

\begin{cousinweb}
\textbf{\sk{सिद्धिः}} (\emph{siddhiḥ}) comes from the root √sidh, \textquotedblleft{}to succeed\textquotedblright{}; MW gives \emph{siddhi} \textquotedblleft{}accomplishment, performance, fulfilment, success\textquotedblright{} (MW).
\end{cousinweb}

\begin{grammarlens}[title={What you've built}]
One more word for saying what happens if.
\end{grammarlens}

\subsection*{Your turn}
\begin{itemize}
\raggedright
  \item \emph{Say it:} \emph{siddhiḥ}
  \item \emph{Say it:} \emph{siddhiḥ}, once more
  \item \emph{Say it:} \emph{yadi prayatnaḥ dṛḍhaḥ, tarhi siddhiḥ avaśyaṃ bhaviṣyati.} — If the effort is firm, success will surely come.
  \item \emph{Recall:} say the Sanskrit for rewards (with a prize), then the Sanskrit for excellence, then say \emph{siddhiḥ} again
\end{itemize}

\begin{tcolorbox}[breakable,colback=teal!4,colframe=teal!35!black,title={Before you move on}]
What does \sk{सिद्धिः} mean? (\textquotedblleft{}An accomplishment\textquotedblright{}.) Say it once more.
\end{tcolorbox}
\section[\texorpdfstring{\sk{सिध्यति} (sidhyati)}{सिध्यति (sidhyati)}]{\sk{सिध्यति} (sidhyati) — succeeds; comes off}

\label{lesson:SA-C386-sidhyati}



\begin{quote}
Before the new one: say the Sanskrit for excellence, then the Sanskrit for an accomplishment.
\end{quote}

\subsection*{You'll want to know: \sk{सिध्यति}}
\textbf{\sk{सिध्यति}} — \emph{sidhyati} — \textquotedblleft{}succeeds; comes off\textquotedblright{}.

\begin{quote}
\sk{यदि} \sk{योजना} \sk{सिध्यति}, \sk{तर्हि} \sk{जनाः} \sk{प्रसन्नाः} \sk{भविष्यन्ति।}
\end{quote}

\begin{quote}
\emph{yadi yojanā sidhyati, tarhi janāḥ prasannāḥ bhaviṣyanti.}
\end{quote}

\begin{quote}
If the plan comes off, people will be pleased.
\end{quote}

Past (it …): \sk{कार्यं} \sk{सिद्धवत्}, \emph{kāryaṃ siddhavat}.

\begin{cousinweb}
\textbf{\sk{सिध्यति}} (\emph{sidhyati}) comes from the root √sidh, \textquotedblleft{}to be accomplished, to succeed\textquotedblright{}; MW gives \emph{siddha} \textquotedblleft{}accomplished, fulfilled, gained\textquotedblright{} (MW).
\end{cousinweb}

\begin{grammarlens}[title={What you've built}]
One more word for saying what happens if.
\end{grammarlens}

\subsection*{Your turn}
\begin{itemize}
\raggedright
  \item \emph{Say it:} \emph{sidhyati}
  \item \emph{Say it:} \emph{sidhyati}, once more
  \item \emph{Say it:} \emph{yadi yojanā sidhyati, tarhi janāḥ prasannāḥ bhaviṣyanti.} — If the plan comes off, people will be pleased.
  \item \emph{Recall:} say the Sanskrit for excellence, then the Sanskrit for an accomplishment, then say \emph{sidhyati} again
\end{itemize}

\begin{tcolorbox}[breakable,colback=teal!4,colframe=teal!35!black,title={Before you move on}]
What does \sk{सिध्यति} mean? (\textquotedblleft{}Succeeds; comes off\textquotedblright{}.) Say it once more.
\end{tcolorbox}
\section[\texorpdfstring{\sk{अर्जयति} (arjayati)}{अर्जयति (arjayati)}]{\sk{अर्जयति} (arjayati) — earns; acquires}

\label{lesson:SA-C386-arjayati}



\begin{quote}
Before the new one: say the Sanskrit for an accomplishment, then the Sanskrit for comes off.
\end{quote}

\subsection*{You'll want to know: \sk{अर्जयति}}
\textbf{\sk{अर्जयति}} — \emph{arjayati} — \textquotedblleft{}earns; acquires\textquotedblright{}.

\begin{quote}
\sk{यदि} \sk{सः} \sk{प्रतिदिनं} \sk{कार्यं} \sk{करोति}, \sk{तर्हि} \sk{धनम्} \sk{अर्जयति।}
\end{quote}

\begin{quote}
\emph{yadi saḥ pratidinaṃ kāryaṃ karoti, tarhi dhanam arjayati.}
\end{quote}

\begin{quote}
If he works every day, he earns money.
\end{quote}

Past (he …): \sk{सः} \sk{बहु} \sk{धनम्} \sk{अर्जितवान्}, \emph{saḥ bahu dhanam arjitavān}.

\begin{cousinweb}
\textbf{\sk{अर्जयति}} (\emph{arjayati}) comes from the root √ṛj, \textquotedblleft{}to procure\textquotedblright{}, in its causative form; MW gives the senses \textquotedblleft{}to procure, to acquire, to obtain\textquotedblright{} and \emph{arjita} \textquotedblleft{}acquired, gained, earned\textquotedblright{} (MW).
\end{cousinweb}

\begin{grammarlens}[title={What you've built}]
One more word for saying what happens if.
\end{grammarlens}

\subsection*{Your turn}
\begin{itemize}
\raggedright
  \item \emph{Say it:} \emph{arjayati}
  \item \emph{Say it:} \emph{arjayati}, once more
  \item \emph{Say it:} \emph{yadi saḥ pratidinaṃ kāryaṃ karoti, tarhi dhanam arjayati.} — If he works every day, he earns money.
  \item \emph{Recall:} say the Sanskrit for an accomplishment, then the Sanskrit for comes off, then say \emph{arjayati} again
\end{itemize}

\begin{tcolorbox}[breakable,colback=teal!4,colframe=teal!35!black,title={Before you move on}]
What does \sk{अर्जयति} mean? (\textquotedblleft{}Earns; acquires\textquotedblright{}.) Say it once more.
\end{tcolorbox}
\section[\texorpdfstring{\sk{अभ्युदयः} (abhyudayaḥ)}{अभ्युदयः (abhyudayaḥ)}]{\sk{अभ्युदयः} (abhyudayaḥ) — a rise in fortune; prosperity}

\label{lesson:SA-C386-abhyudayah}



\begin{quote}
Before the new one: say the Sanskrit for comes off, then the Sanskrit for earns.
\end{quote}

\subsection*{You'll want to know: \sk{अभ्युदयः}}
\textbf{\sk{अभ्युदयः}} — \emph{abhyudayaḥ} — \textquotedblleft{}a rise in fortune; prosperity\textquotedblright{}.

\begin{quote}
\sk{यदि} \sk{शान्तिः} \sk{अस्ति}, \sk{तर्हि} \sk{देशस्य} \sk{अभ्युदयः} \sk{भविष्यति।}
\end{quote}

\begin{quote}
\emph{yadi śāntiḥ asti, tarhi deśasya abhyudayaḥ bhaviṣyati.}
\end{quote}

\begin{quote}
If there is peace, the country will prosper.
\end{quote}

A masculine noun.

\begin{cousinweb}
\textbf{\sk{अभ्युदयः}} (\emph{abhyudayaḥ}) is built on \emph{abhi}, \emph{ud} and the root √i, \textquotedblleft{}to go\textquotedblright{}; MW gives \emph{abhy-udaya} \textquotedblleft{}sunrise\textquotedblright{} and \textquotedblleft{}elevation, increase, prosperity\textquotedblright{} (MW).
\end{cousinweb}

\begin{grammarlens}[title={What you've built}]
One more word for saying what happens if.
\end{grammarlens}

\subsection*{Your turn}
\begin{itemize}
\raggedright
  \item \emph{Say it:} \emph{abhyudayaḥ}
  \item \emph{Say it:} \emph{abhyudayaḥ}, once more
  \item \emph{Say it:} \emph{yadi śāntiḥ asti, tarhi deśasya abhyudayaḥ bhaviṣyati.} — If there is peace, the country will prosper.
  \item \emph{Recall:} say the Sanskrit for comes off, then the Sanskrit for earns, then say \emph{abhyudayaḥ} again
\end{itemize}

\begin{tcolorbox}[breakable,colback=teal!4,colframe=teal!35!black,title={Before you move on}]
What does \sk{अभ्युदयः} mean? (\textquotedblleft{}A rise in fortune; prosperity\textquotedblright{}.) Say it once more.
\end{tcolorbox}
\section[\texorpdfstring{\sk{समृद्धिः} (samṛddhiḥ)}{समृद्धिः (samṛddhiḥ)}]{\sk{समृद्धिः} (samṛddhiḥ) — prosperity}

\label{lesson:SA-C386-samrddhih}



\begin{quote}
Before the new one: say the Sanskrit for earns, then the Sanskrit for a rise in fortune.
\end{quote}

\subsection*{You'll want to know: \sk{समृद्धिः}}
\textbf{\sk{समृद्धिः}} — \emph{samṛddhiḥ} — \textquotedblleft{}prosperity\textquotedblright{}.

\begin{quote}
\sk{यदि} \sk{व्यापारः} \sk{वर्धते}, \sk{तर्हि} \sk{नगरे} \sk{समृद्धिः} \sk{भविष्यति।}
\end{quote}

\begin{quote}
\emph{yadi vyāpāraḥ vardhate, tarhi nagare samṛddhiḥ bhaviṣyati.}
\end{quote}

\begin{quote}
If trade grows, there will be prosperity in the city.
\end{quote}

A feminine noun.

\begin{cousinweb}
\textbf{\sk{समृद्धिः}} (\emph{samṛddhiḥ}) is built on \emph{sam} and the root √ṛdh, \textquotedblleft{}to thrive\textquotedblright{}; MW gives \emph{sam-ṛddhi} \textquotedblleft{}great prosperity or success, growth, welfare\textquotedblright{} (MW).
\end{cousinweb}

\begin{grammarlens}[title={What you've built}]
One more word for saying what happens if.
\end{grammarlens}

\subsection*{Your turn}
\begin{itemize}
\raggedright
  \item \emph{Say it:} \emph{samṛddhiḥ}
  \item \emph{Say it:} \emph{samṛddhiḥ}, once more
  \item \emph{Say it:} \emph{yadi vyāpāraḥ vardhate, tarhi nagare samṛddhiḥ bhaviṣyati.} — If trade grows, there will be prosperity in the city.
  \item \emph{Recall:} say the Sanskrit for earns, then the Sanskrit for a rise in fortune, then say \emph{samṛddhiḥ} again
\end{itemize}

\begin{tcolorbox}[breakable,colback=teal!4,colframe=teal!35!black,title={Before you move on}]
What does \sk{समृद्धिः} mean? (\textquotedblleft{}Prosperity\textquotedblright{}.) Say it once more.
\end{tcolorbox}
\section[(dialogue)]{Looking back: from How Much to Success}

\label{lesson:SA-R386-review-380-386}



\begin{quote}
Say the Sanskrit for a quantity and an amount.
\end{quote}

\subsection*{Your turn}
Say these aloud:

\begin{itemize}
\raggedright
  \item a quantity, an amount, moderate, as much as one likes, expands
  \item generally, especially, in every way, repeatedly, sums up
  \item benefit, welfare, the better course, rests, receives with honour
  \item a resolve, an ambition, longs for, a heartfelt wish, more and more
  \item perseverance, constancy, as far as one can, carries out, encourages
  \item competes, a rival, a winner, rewards (with a prize), excellence
  \item an accomplishment, comes off, earns, a rise in fortune, prosperity
\end{itemize}

\begin{tcolorbox}[breakable,colback=teal!4,colframe=teal!35!black,title={Before you move on}]
A quantity? (\textbf{\sk{मात्रा}}, \emph{mātrā}.) An amount? (\textbf{\sk{परिमाणम्}}, \emph{parimāṇam}.) Moderate? (\textbf{\sk{मध्यमः}}, \emph{madhyamaḥ}.) As much as one likes? (\textbf{\sk{यथेष्टम्}}, \emph{yatheṣṭam}.) Expands? (\textbf{\sk{विस्तारयति}}, \emph{vistārayati}.) Generally? (\textbf{\sk{सामान्यतः}}, \emph{sāmānyataḥ}.) Especially? (\textbf{\sk{विशेषतः}}, \emph{viśeṣataḥ}.) In every way? (\textbf{\sk{सर्वथा}}, \emph{sarvathā}.) Repeatedly? (\textbf{\sk{बहुशः}}, \emph{bahuśaḥ}.) Sums up? (\textbf{\sk{संक्षिपति}}, \emph{saṅkṣipati}.) Benefit? (\textbf{\sk{हितम्}}, \emph{hitam}.) Welfare? (\textbf{\sk{कल्याणम्}}, \emph{kalyāṇam}.) The better course? (\textbf{\sk{श्रेयः}}, \emph{śreyaḥ}.) Rests? (\textbf{\sk{विश्राम्यति}}, \emph{viśrāmyati}.) Receives with honour? (\textbf{\sk{सत्करोति}}, \emph{satkaroti}.) A resolve? (\textbf{\sk{संकल्पः}}, \emph{saṅkalpaḥ}.) An ambition? (\textbf{\sk{आकांक्षा}}, \emph{ākāṅkṣā}.) Longs for? (\textbf{\sk{आकांक्षति}}, \emph{ākāṅkṣati}.) A heartfelt wish? (\textbf{\sk{अभिलाषः}}, \emph{abhilāṣaḥ}.) More and more? (\textbf{\sk{उत्तरोत्तरम्}}, \emph{uttarottaram}.) Perseverance? (\textbf{\sk{उद्यमः}}, \emph{udyamaḥ}.) Constancy? (\textbf{\sk{धृतिः}}, \emph{dhṛtiḥ}.) As far as one can? (\textbf{\sk{यथाशक्ति}}, \emph{yathāśakti}.) Carries out? (\textbf{\sk{अनुतिष्ठति}}, \emph{anutiṣṭhati}.) Encourages? (\textbf{\sk{प्रोत्साहयति}}, \emph{protsāhayati}.) Competes? (\textbf{\sk{स्पर्धते}}, \emph{spardhate}.) A rival? (\textbf{\sk{प्रतिद्वन्द्वी}}, \emph{pratidvandvī}.) A winner? (\textbf{\sk{जेता}}, \emph{jetā}.) Rewards (with a prize)? (\textbf{\sk{पुरस्करोति}}, \emph{puraskaroti}.) Excellence? (\textbf{\sk{उत्कर्षः}}, \emph{utkarṣaḥ}.) An accomplishment? (\textbf{\sk{सिद्धिः}}, \emph{siddhiḥ}.) Comes off? (\textbf{\sk{सिध्यति}}, \emph{sidhyati}.) Earns? (\textbf{\sk{अर्जयति}}, \emph{arjayati}.) A rise in fortune? (\textbf{\sk{अभ्युदयः}}, \emph{abhyudayaḥ}.) Prosperity? (\textbf{\sk{समृद्धिः}}, \emph{samṛddhiḥ}.)
\end{tcolorbox}
